%% kapa: exam=1 round=34 year=2023 period=1 form=A subject=real_estate subject_ko=부동산학원론 seq=32 number=112
%% source: Q-Net 감정평가사 기출문제 (cst00309) artlSeq=5211789 file=raw/1008/5211789/2023년도 제34회 감정평가사 1차시험 1교시.pdf
%% conversion: kordoc 4.24.51 -> tools/gen1.py (auto); review=proofread-sonnet; flags=none
\begin{question}{112}
우리나라의 부동산등기제도와 권리분석에 관한 설명으로 옳지 않은 것은?
\begin{choices}
  \item 소유권이전등기 청구권을 확보하기 위해 처분금지가처분의 등기가 가능하다.
  \item 현재 환매(특약)등기제와 예고등기제는 「부동산등기법」상 폐지되었다.
  \item 등기의 순서는 같은 구(區)에서 한 등기 상호간에는 순위번호에 따른다.
  \item 근저당권과 담보가등기는 부동산경매에서 말소기준권리가 된다.
  \item 부동산임차권은 부동산물권이 아니지만 등기할 수 있다.
\end{choices}
\end{question}
\answer{2}
